\documentclass[10pt,a4paper]{amsart}
\usepackage[margin=19mm]{geometry}
\usepackage{amsmath,amssymb,mathtools}
\usepackage[scr=rsfs,cal=euler]{mathalfa}
\usepackage{xfrac}
\usepackage[unicode,hidelinks,bookmarks=false]{hyperref}
\newcommand{\ie}{i.e.\ }
\newcommand{\cf}{cf.\ }
\newcommand{\resp}{respectively\ }
\newcommand{\Subsection}{\subsection}
\providecommand{\nobreakdash}{\nobreak\hbox{-}}
\setlength{\emergencystretch}{2em}
\multlinegap0pt
\usepackage{bm}
\usepackage{bbm}
\usepackage{dsfont}
\usepackage{xspace}
\usepackage{cancel}
\usepackage{mathtools}
\usepackage{amsthm}
\makeatletter
\@ifundefined{theorem}{\newtheorem{theorem}{Theorem}}{}
\@ifundefined{lemma}{\newtheorem{lemma}[theorem]{Lemma}}{}
\makeatother
\newcommand{\dt} {\displaystyle\mathrm{d}t}
\newcommand{\dx} {\displaystyle\mathrm{d}x}
\title[Reconstruction of a potential parameter in subdiffusion via ]{Reconstruction of a potential parameter in subdiffusion via a Kohn--Vogelius type functional: Theory and computation}
\author{Hamza Kahlaoui, Mourad Hrizi, Abdessamad Oulmelk, Xiangcheng Zheng, Mahmoud A. Zaky,, Ahmed Hendy}
\date{Source: arXiv:2509.19260v2 (CC BY 4.0)}
\begin{document}
\maketitle

\noindent\emph{Source and licence.} Reconstruction of a potential parameter in subdiffusion via a Kohn--Vogelius type functional: Theory and computation. Version arXiv:2509.19260v2 (CC BY 4.0), distributed under the Creative Commons Attribution 4.0 International licence, as recorded in the arXiv licence metadata for this exact version and shown on the abstract page https://arxiv.org/abs/2509.19260v2. Full rights notice: licenses/arxiv-2509.19260-CC-BY-4.0.txt.\par\smallskip

\noindent\textit{Editorial scope.} Counted unit: the Theorem whose source label is \texttt{None} in the article (source0.tex lines 1195--1200, source label \texttt{None}), delivered together with the article's own complete original proof of it (source0.tex lines 1202--1314, one \texttt{proof} environment of 113 lines). No mathematics is written, repaired or re-derived: statement and proof are the article's own text line for line. Required context retained uncounted: PN (lines 418--426); estimate-Direchlet (lines 440--442); stable (lines 463--466); h1+ (lines 526--528); h11 (lines 534--537); conv\_N (lines 560--566); theorem-existence (lines 577--580); cv-0 (lines 591--593); c2 (lines 648--652); new-uniqueness-condition (lines 1023--1025); lim-seq-sat (lines 1164--1166). Every definition, display and label of those lines is retained unchanged. Omitted from this package and not redistributed: 1--417, 427--439, 443--462, 467--525, 529--533, 538--559, 567--576, 581--590, 594--647, 653--1022, 1026--1163, 1167--1194, 1201--1201, 1315--2708. Nothing inside a retained excerpt is omitted or altered: every retained body is delivered exactly as the article's lines are written, so no omission receipt of any kind accompanies this package. Citations: the retained excerpts carry no citation command at all, so no bibliography entry of the article is transcribed and no reference is redistributed. Reference adaptation: none was needed and none was made; every reference inside the delivered unit resolves inside it, so no unresolved reference or citation is produced. Printed slips kept literal: no printed word, symbol, hypothesis or number of the counted statement or of its proof was altered. Layout and class: the article's own class is \texttt{amsart} with packages including amsmath, amssymb, amsthm, dsfont, stmaryrd, float, color, tabu, graphics, graphicx, subfigure, booktabs, multirow, array; the wrapper is the lane's pinned \texttt{amsart} editorial wrapper at 10pt on A4, which restates the theorem-like environments the retained text needs and adds the article's own macro definitions that the retained text needs (\texttt{\textbackslash dt}, \texttt{\textbackslash dx}), with extra wrapper packages (bm, bbm, dsfont, xspace, cancel, mathtools). The delivered numbering is the unit's own. Assets: no retained span contains a figure environment, a graphics-inclusion command, a PDF-page inclusion, a table or a TikZ picture; no image file is copied into this package, the package declares no asset dependency and no image file is redistributed with it. Licence: version arXiv:2509.19260v2 of this article is distributed under the Creative Commons Attribution 4.0 International licence, as recorded in the arXiv licence metadata for this exact version and shown on the abstract page https://arxiv.org/abs/2509.19260v2. A full rights notice is delivered at licenses/arxiv-2509.19260-CC-BY-4.0.txt. Characters: every retained excerpt is pure ASCII, so the delivered engine, XeLaTeX with its default Latin Modern text font, reports no missing character.
\begin{eqnarray} \label{PN}
\left\{
  \begin{array}{rcll}
    \partial_{t} ^{\alpha}  u_N[q]-\Delta  u_N[q] +q\, u_N[q]&=&0&\hbox{in}\ Q_T, \\
    \displaystyle \partial_\nu u_N[q]&=&\varphi &\hbox{on}\  \Sigma_T, \\
   u_N[q](\cdot,0)&=&0  &\hbox{in}\ \Omega.
  \end{array}
\right.
\end{eqnarray}

\begin{equation} \label{estimate-Direchlet}
    \| u_D[q] \|_{B^{\frac{\alpha}{2}}(Q_T)} \leq C_D \|\phi\|_{H^{\frac{1}{2}, \frac{\alpha}{4}}(\Sigma_T)},
\end{equation}

\begin{equation}\label{stable} 
    \displaystyle\operatorname*{Minimize}_{q\in\Phi_{ad}}\
\mathcal{K}_\rho( q):=\mathcal{K}( q)+\rho\|q\|^2_{L^2(\Omega)},
\end{equation}

\begin{align} \label{h1+}
    \|w_n\|^2_{L^2(0,T;H^1(\Omega))} \leq C \int_0^T\int_{\Omega}\Big|(q^n -q) u_D[q]w_n\Big| \dx \dt,
\end{align}

\begin{align} \label{h11}
\int_0^T\int_\Omega\Big|(q^n &-q) u_D[q]\, w_n\Big|\dx\dt\nonumber\\&\leq C(\Omega,d)
\|q^n -q\|_{L^2(\Omega)}\| u_D[q]\|_{L^2(0,T;H^1(\Omega))}\|w_n\|_{L^2(0,T;H^1(\Omega))}.
\end{align}

\begin{lemma} \label{conv_N}
Let $q\in\Phi_{ad}$ and $q^n$ be a sequence of functions in $\Phi_{ad}.$ If $\{ q^n \}_n$ converges strongly to  $q$ in $L^2(\Omega)$. Then
\begin{equation} \label{convergancefaible-N}
    u_N[q^n]\longrightarrow u_N[q]\; \hbox{ strongly in }\; L^2(0,T;H^1(\Omega))\; \hbox{ as }\; n\to\infty, 
\end{equation}
where $u_N[q]$ is the solution of \eqref{PN} and $u_N[q^n]$ is the solution of \eqref{PN} with $q^n$ replacing $q$.
\end{lemma}

\begin{theorem}[Existence]\label{theorem-existence}
For every $\rho>0$ there exists at least one solution of the optimization
problem \eqref{stable}.
\end{theorem}

\begin{equation}\label{cv-0}
q^n\longrightarrow q^0\; \hbox{strongly in}\; L^2(\Omega)\; \hbox{as}\; n\rightarrow\infty.
\end{equation}

\begin{align} \label{c2}
    \lim_{n\rightarrow\infty} \int_0^T \int_{\Omega}q^n\Big|
 u_N[q^n]- u_D[q^n]\Big|^2\ \dx \dt = \int_0^T \int_{\Omega}q^0\Big|
 u_N[q^0]- u_D[q^0]\Big|^2\ \dx \dt.
\end{align}

\begin{theorem}[Uniqueness]\label{new-uniqueness-condition}
    For any $\rho>0$, there exists a positive constant $M$ independent of $q\in\Phi_{ad}$ such that for $M<\rho$, $\mathcal{K}_\rho(q)$ is strictly convex.
\end{theorem}

\begin{align}\label{lim-seq-sat}
    \|\phi^n - \phi\|_{H^{\frac{1}{2},\frac{\alpha}{4} }\left(\Sigma_{T}\right)} \longrightarrow 0\;\; \hbox{as}\; n\rightarrow\infty.
\end{align}

\begin{theorem} Let $\rho>0$ satisfy the uniqueness condition in Theorem \ref{new-uniqueness-condition} and let $\{\phi^n\} \subset {}_0H^{\frac{1}{2},\frac{\alpha}{4}}(\Sigma_T)$ be a sequence of perturbed measurements satisfying \eqref{lim-seq-sat}. Let $\{ q^n \}\subset\Phi_{ad}$ be a sequence of minimizers of problems $(\mathcal{O}\mathcal{P}_{\rho}^n)$ for $n=1,2,\cdots$. Then there exists a subsequence $\{ q^{n_\ell} \}$ of $\{ q^n \}$, such that 
\begin{equation*}
q^{n_\ell}\longrightarrow q^0\; \hbox{strongly in}\; L^2(\Omega)\; \hbox{as}\; \ell\rightarrow\infty,
\end{equation*}
where $q^0\in\mathcal{Q}$ is the unique minimizer of the optimization problem \eqref{stable}.
\end{theorem}

\begin{proof}
The unique existence of each $q^n\in\Phi_{ad}$ is ensured by Theorems \ref{theorem-existence} and \ref{new-uniqueness-condition}. Employing a similar technique as in the proof of the convergence result \eqref{cv-0}, we establish that the sequence $q^n$ strongly converges to $q^0\in\Phi_{ad}$ in $L^2(\Omega)$. In other words, we have
\begin{equation}\label{cvv-st}
q^n\longrightarrow q^0\; \hbox{strongly in}\; L^2(\Omega)\; \hbox{as}\; n\rightarrow\infty.
\end{equation}


Now it suffices to show that $q^0$ is indeed the unique minimizer of \eqref{stable}. From \eqref{cvv-st} and Lemma \ref{conv_N}, one can deduce that 
\begin{equation} \label{convergancefaible-NNN}
    u_N[q^n]\longrightarrow u_N[q^0]\; \hbox{ strongly in }\; L^2(0,T;H^1(\Omega))\; \hbox{ as }\; n\to\infty, 
\end{equation}
where $u_N[q^0]$ is the solution to \eqref{PN} with $q=q^0$, and $u_N[q^n]$ is the solution to \eqref{PN} with $q=q^n$.

Next, we prove that
\begin{equation}\label{convergancefaible-DDD}
    u_D[q^n, \phi^n] \longrightarrow u_D[q^0, \phi] \quad \text{stongly in} \quad L^2(0, T; L^2(\Omega)) \;\; \text{as} \;\; n \longrightarrow \infty.
\end{equation}
To establish the above convergence result, we define $\vartheta_n = u_D[q^n, \phi^n] - u_D[q^0, \phi]$. It is straightforward to deduce that $\vartheta_n$ is the solution to
\begin{eqnarray*} %\label{P_d22}
\left\{
  \begin{array}{rcll}
    \partial_t^\alpha \vartheta_n - \Delta \vartheta_n + q^0 \vartheta_n &=& -(q^n - q^0) u_D[q^n, \phi^n] & \text{in} \; Q_T, \\
    \vartheta_n &=& \phi^n - \phi & \text{on} \; \Sigma_T, \\
    \vartheta_n(\cdot, 0) &=& 0 & \text{in} \; \Omega.
  \end{array}
\right.
\end{eqnarray*}
We split $\vartheta_n$ into $\vartheta_n^1 + \vartheta_n^2$ such that
\begin{eqnarray} \label{P_d30}
\left\{
  \begin{array}{rcll}
    \partial_t^\alpha \vartheta_n^1 - \Delta \vartheta_n^1 + q^0 \vartheta_n^1 &=& -(q^n - q^0) u_D[q^n, \phi^n] & \text{in} \; Q_T, \\
    \vartheta_n^1 &=& 0 & \text{on} \; \Sigma_T, \\
    \vartheta_n^1(\cdot, 0) &=& 0 & \text{in} \; \Omega.
  \end{array}
\right.
\end{eqnarray}
\begin{eqnarray*} %\label{P_d3}
\left\{
  \begin{array}{rcll}
    \partial_t^\alpha \vartheta_n^2 - \Delta \vartheta_n^2 + q^0 \vartheta_n^2 &=& 0 & \text{in} \; Q_T, \\
    \vartheta_n^2 &=& \phi^n - \phi & \text{on} \; \Sigma_T, \\
    \vartheta_n^2(\cdot, 0) &=& 0 & \text{in} \; \Omega.
  \end{array}
\right.
\end{eqnarray*}
By taking $\vartheta_n^1$ as a test function in the weak formulation of \eqref{P_d30}, we get
\begin{align*}
    \int_0^T \int_\Omega \Big(\partial_t^\alpha \vartheta_n^1\Big) \vartheta_n^1 \dx \dt + \int_0^T \int_\Omega \Big|\nabla \vartheta_n^1\Big|^2 \dx \dt + \int_0^T \int_\Omega q^0 \Big|\vartheta_n^1\Big|^2 \dx \dt \\
    = \int_0^T \int_\Omega \Big(q^n - q^0\Big) \vartheta_n^1 u_D[q^n, \phi^n] \, \dx \dt.
\end{align*}
Employing a similar technique as in the proof of inequality \eqref{h1+}, we obtain
\begin{equation} \label{dd_1}
    \|\vartheta_n^1\|_{L^2(0, T; H^1(\Omega))}^2 \leq C(\Omega) \int_0^T \int_\Omega \Big|\Big(q^n - q^0\Big) \vartheta_n^1 u_D[q^n, \phi^n]\Big| \dx \dt.
\end{equation}
Similar to \eqref{h11}, one can deduce
\begin{equation}\label{dd-ff}
    \int_0^T \int_\Omega \Big|\Big(q^n - q^0\Big) \vartheta_n^1 u_D[q^n, \phi^n]\Big| \dx \dt \leq C(\Omega) \|q^n - q^0\|_{L^2(\Omega)} \|\vartheta_n^1\|_{L^2(0, T; H^1(\Omega))} \| u_D[q^n, \phi^n] \|_{L^2(0, T; H^1(\Omega))}.
\end{equation}
From the inequality \eqref{estimate-Direchlet}, we obtain
\begin{align*}
    \| u_D[q^n, \phi^n] \|_{L^2(0, T; H^1(\Omega))} &\leq C(\alpha, T, \Omega) \|\phi^n\|_{H^{\frac{1}{2}, \frac{\alpha}{4}} (\Sigma_T)} \\
    &\leq C(\alpha, T, \Omega) \Big(\|\phi^n - \phi\|_{H^{\frac{1}{2}, \frac{\alpha}{4}} (\Sigma_T)} + \|\phi\|_{H^{\frac{1}{2}, \frac{\alpha}{4}} (\Sigma_T)}\Big).
\end{align*}
Using the fact that \(\phi^n \to \phi\) in \({}_0 H^{\frac{1}{2}, \frac{\alpha}{4}} (\Sigma_T)\) as \(n \to \infty\), we obtain
\begin{align}\label{insert}
    \| u_D[q^n, \phi^n] \|_{L^2(0, T; H^1(\Omega))} \leq C(\alpha, T, \Omega, \phi).
\end{align}
Inserting \eqref{insert} into \eqref{dd-ff}, we obtain
\begin{equation}\label{dd-ff2}
    \int_0^T \int_\Omega \Big|\Big(q^n - q^0\Big) \vartheta_n^1 u_D[q^n, \phi^n]\Big| \dx \dt \leq C(\alpha, T, \Omega, \phi) \|q^n - q^0\|_{L^2(\Omega)} \|\vartheta_n^1\|_{L^2(0, T; H^1(\Omega))}.
\end{equation}
Inserting \eqref{dd-ff2} into \eqref{dd_1}, we get
\begin{equation}\label{dd-ff3}
    \|\vartheta_n^1\|_{L^2(0, T; H^1(\Omega))} \leq C(\alpha, T, \Omega,\phi) \|q^n - q^0\|_{L^2(\Omega)}.
\end{equation}
Moreover, from the inequality \eqref{estimate-Direchlet}, we obtain
\begin{align} \label{dd_ff4}
    \|\vartheta_n^2\|_{L^2(0, T; H^1(\Omega))} \leq C(\alpha, T, \Omega) \|\phi^n - \phi\|_{L^2(0, T; L^2(\partial\Omega))}.
\end{align}
Combining \eqref{dd-ff3} with \eqref{dd_ff4}, and using the convergence result \eqref{cvv-st} and the fact that $\phi^n \longrightarrow \phi$ in ${}_0 H^{\frac{1}{2}, \frac{\alpha}{4}} (\Sigma_T)$ as $n \longrightarrow \infty$, we get
\begin{equation*}
    \|\vartheta_n\|_{L^2(0, T; H^1(\Omega))} \longrightarrow 0 \;\; \text{as} \;\; n \longrightarrow \infty.
\end{equation*}

On the other hand, following the same strategy developed in the proof of \eqref{c2}, we obtain
\begin{align} \label{c2++}
    \lim_{n \rightarrow \infty} \int_0^T \int_\Omega q^n \Big| u_N[q^n] - u_D[q^n, \phi^n] \Big|^2 \dx \dt = \int_0^T \int_\Omega q^0 \Big| u_N[q^0] - u_D[q^0, \phi] \Big|^2 \dx \dt.
\end{align}

Consequently, by combining the convergence results \eqref{cvv-st}, \eqref{convergancefaible-NNN}, \eqref{convergancefaible-DDD}, and \eqref{c2++} along with the lower semi-continuity of the $L^2$-norm, we conclude

\begin{align*}
\mathcal{K}_\rho(q^0) &=\int_0^T \int_{\Omega}\Big|\nabla\Big(
 u_N[q^0]- u_D[q^0]\Big)\Big|^2\ \dx \dt + \int_0^T \int_{\Omega}q^0\Big|
 u_N[q^0]- u_D[q^0]\Big|^2\ \dx\dt\\
&\;\;\quad +\rho\int_0^T \int_{\Omega}\Big|q^0\Big|^2\ \dx\dt\\
&\leq\lim_{n\rightarrow\infty}\inf\int_0^T \int_{\Omega}\Big|\nabla\Big(
 u_N[q^n]- u_D[q^n,\phi^n]\Big)\Big|^2\ \dx \dt \\
&\;\;\quad + \lim_{n\rightarrow\infty}\inf\int_0^T \int_{\Omega}q^n\Big|
 u_N[q^n]- u_D[q^n,\phi^n]\Big|^2\ \dx\dt+ \rho\lim_{n\rightarrow\infty}\inf\int_0^T \int_{\Omega}\Big|q^n\Big|^2\ \dx\dt\\
&\leq \lim_{n\rightarrow\infty}\inf \mathcal{K}^n_\rho(q^n).
\end{align*}
Moreover, since \(\{ q^n \} \subset \Phi_{ad}\) is a sequence of minimizers of problems \((\mathcal{O}\mathcal{P}_{\rho}^n)\), and by using \(u_D[q,\phi^n] \longrightarrow u_D[q,\phi]\) in \(L^2(0, T; H^1(\Omega))\) as \(n \longrightarrow \infty\), we obtain
\begin{align*} 
\mathcal{K}_\rho(q^0)&\leq\lim_{n\rightarrow\infty}\Big(\int_0^T \int_{\Omega}\Big|\nabla\Big(
 u_N[q]- u_D[q,\phi^n]\Big)\Big|^2\ \dx \dt + \int_0^T \int_{\Omega}q\Big|
 u_N[q]- u_D[q,\phi^n]\Big|^2\ \dx\dt \\
 &\;\; \quad + \rho\int_0^T \int_{\Omega}\Big|q\Big|^2\ \dx\dt \Big)\\
&=\mathcal{K}_\rho(q), \quad \;\;\; \hbox{for any} \; q \in \Phi_{ad},
\end{align*}
which verifies that $q^0$ is the minimizer of \eqref{stable}.
\end{proof}

\end{document}
